\documentclass[11pt]{article}
\usepackage[margin=1in]{geometry}
\usepackage{amsmath,amssymb}
\usepackage[hidelinks]{hyperref}

\title{Response: the CSL/SCG open issue is not open}
\date{}

\begin{document}
\maketitle
\thispagestyle{empty}

\noindent
The note is a faithful rendering of both what was said in discussion and
what \texttt{PROOFS\_v2.tex} currently states. Both are out of date on
this point, and the note inherits the staleness in a way that matters:
it presents as an unresolved question something that has already been
computed, and it proposes a route to closing it that is already known to
fail. What follows corrects each in turn and states what should replace
them.

\bigskip
\noindent\textbf{The estimated ratio has been computed.}
The note says that whether the estimated variance ratio is invariant to
the preference parameter ``has not been established.'' It has.
\texttt{lambda\_RN\_measured.py} simulates the solved model, samples it
quarterly over a sixty-quarter window after a long burn-in, and applies
the panel's own estimator---the same estimator that would be used on real
data---at four hundred paths per value of the asymmetry parameter. Under
the first cost convention the median estimated ratio runs from
approximately $1.334$ at the neutral parameter value down to
approximately $1.152$ at the top of the admissible range; under the second
it runs from approximately $1.390$ to $1.167$. The estimated ratio
therefore \emph{does} vary with the preference parameter, it decreases,
and the direction is known.

The sentence in \texttt{PROOFS\_v2.tex} asserting that this ``is not
known'' was already flagged for revision before this note was written. It
was never edited, and the note has now carried it into a second document.
That is precisely how a stale line becomes a stale literature, and it is
the strongest practical argument for making the edit rather than
continuing to work around it.

\bigskip
\noindent\textbf{The proposed closure route is already known to fail.}
The note proposes closing the gap by showing that the estimated ratio
``converges to or otherwise tracks the definitional one as sample size
grows.'' That cannot be shown, because it is false. The panel estimator is
asymptotically biased toward one: a true value of $1.30$ measures $1.164$
at sixty observations, $1.145$ at five hundred, and $1.140$ at twenty
thousand. The bias converges to roughly $-0.16$ and stays there. This is
not a finite-sample deficiency that more data repairs; the estimator
targets a \emph{different estimand}, because states near the
classification boundary are assigned to one regime while their innovations
reflect a mixture of both. The bridge the note asks for does not exist and
will not be built. Treating this as pending an asymptotic argument points
at a question that has already returned a negative answer.

\bigskip
\noindent\textbf{What should replace the open-issue framing.}
The correct statement is not that volatility is uninformative and the
estimated case is unchecked. It is a pair of claims at different
strengths. First, the definitional evaluation of the ratio carries no
information about the preference whatsoever---this is proved (SCG), and it
is a derived consequence of how the cap's geometry propagates into the
saturation limits, not a definitional identity. Second, the estimated
ratio carries a weak signal, and the weakness is quantifiable: roughly
$0.18$ of total movement across the entire admissible range of the
preference parameter, against a per-path standard deviation of $0.16$ to
$0.20$. For a single institution the signal is indistinguishable from its
own noise; only a median across many institutions is informative at all.

That is a quantitative weak-identification claim rather than a blank
uninformativeness claim, and it is the stronger position of the two. It
survives the obvious objection---someone will estimate the ratio, obtain a
number that moves with the parameter, and the paper needs to have
predicted that movement and characterised its size rather than been
surprised by it.

\bigskip
\noindent\textbf{Consequence for the choice of first empirical test.}
The note's conclusion is right and its reasoning is not. The K-sweep
prediction (KSW) is the correct first test, but not because the
definitional-versus-estimated relationship is unresolved groundwork. It is
correct because that relationship \emph{is} resolved, and unfavourably:
the estimated ratio is a biased, noisy, weakly-responsive proxy for a
quantity that is provably constant. A threshold-based test therefore
dominates a variance-based one on identification grounds, which is a
positive argument for KSW rather than a deferral of the alternative.

\bigskip
\noindent\textbf{Two smaller points.}
The phrase ``true by construction rather than by measurement'' undersells
SCG and invites a referee to read it as tautological; it is a proved
symbolic result, and ``proved from the model's structure'' carries the
same meaning without the concession. And the suggestion that closure would
proceed ``presumably through simulation'' describes work already performed
at four hundred paths per parameter value, which will read oddly to anyone
who has seen the measured-scale computation.

\bigskip
\noindent\textbf{Shortest path to fixing this.}
One edit to the open-questions section of \texttt{PROOFS\_v2.tex},
replacing the ``is not known'' sentence with the measured direction and
the weak-identification magnitude, followed by a rewrite of this note
against the corrected text. The note's structure survives; only its
premise changes.

\end{document}
